
\begin{figure}[H]
    \centering
    \begin{tikzpicture}[
        font=\small,
        box/.style={draw, rounded corners, align=center, text width=0.29\linewidth, minimum height=1.1cm, fill=gray!8},
        data/.style={draw, rounded corners, align=center, text width=0.24\linewidth, minimum height=0.9cm, fill=blue!7},
        >=Latex,
        node distance=7mm and 8mm,
    ]
        \node[box] (rq1) {RQ1\\Parametric retrofit adaptability};
        \node[box, right=of rq1] (rq2) {RQ2\\Low-cost motion and positioning};
        \node[box, right=of rq2] (rq3) {RQ3\\Calibration-enabled workflows};
        \node[data, below=of rq1] (cad) {CAD parameter variants\\Mechanical artifact};
        \node[data, below=of rq2] (motion) {Stage scale\\Motion accuracy\\Calibration repeat};
        \node[data, below=of rq3] (workflow) {Autofocus\\Focus stacking\\Tile scanning};
        \draw[-{Latex}] (cad) -- (rq1);
        \draw[-{Latex}] (motion) -- (rq2);
        \draw[-{Latex}] (motion) -- (rq3);
        \draw[-{Latex}] (workflow) -- (rq3);
    \end{tikzpicture}
    \caption[Evaluation evidence map]{Evaluation evidence map. The physical measurements mainly support RQ2 and RQ3, while RQ1 is supported by the mechanical artifact and virtual parametric adaptation rather than by a second physical microscope build.}
    \label{fig:eval-evidence-map}
\end{figure}
